\documentclass[10pt,a4paper]{amsart}
\usepackage{amsmath,amssymb,amsthm,graphicx}
\usepackage{verbatim}
\newenvironment{reference}{\comment}{\endcomment}
%\newenvironment{reference}{}{}
\newenvironment{slogan}{\comment}{\endcomment}
\newenvironment{history}{\comment}{\endcomment}
% For commutative diagrams we use Xy-pic
\usepackage[all]{xy}

% We use 2cell for 2-commutative diagrams.
\xyoption{2cell}
% Theorem environments.
%
\theoremstyle{plain}
\newtheorem{theorem}[subsection]{Theorem}
\newtheorem{proposition}[subsection]{Proposition}
\newtheorem{lemma}[subsection]{Lemma}

\theoremstyle{definition}
\newtheorem{definition}[subsection]{Definition}
\newtheorem{example}[subsection]{Example}
\newtheorem{exercise}[subsection]{Exercise}
\newtheorem{situation}[subsection]{Situation}

\theoremstyle{remark}
\newtheorem{remark}[subsection]{Remark}
\newtheorem{remarks}[subsection]{Remarks}

\numberwithin{equation}{subsection}

% Macros
%
\def\lim{\mathop{\mathrm{lim}}\nolimits}
\def\colim{\mathop{\mathrm{colim}}\nolimits}
\def\Spec{\mathop{\mathrm{Spec}}}
\def\Hom{\mathop{\mathrm{Hom}}\nolimits}
\def\Ext{\mathop{\mathrm{Ext}}\nolimits}
\def\SheafHom{\mathop{\mathcal{H}\!\mathit{om}}\nolimits}
\def\SheafExt{\mathop{\mathcal{E}\!\mathit{xt}}\nolimits}
\def\Sch{\mathit{Sch}}
\def\Mor{\mathop{\mathrm{Mor}}\nolimits}
\def\Ob{\mathop{\mathrm{Ob}}\nolimits}
\def\Sh{\mathop{\mathit{Sh}}\nolimits}
\def\NL{\mathop{N\!L}\nolimits}
\def\CH{\mathop{\mathrm{CH}}\nolimits}
\def\proetale{{pro\text{-}\acute{e}tale}}
\def\etale{{\acute{e}tale}}
\def\QCoh{\mathit{QCoh}}
\def\Ker{\mathop{\mathrm{Ker}}}
\def\Im{\mathop{\mathrm{Im}}}
\def\Coker{\mathop{\mathrm{Coker}}}
\def\Coim{\mathop{\mathrm{Coim}}}

% Boxtimes
%
\DeclareMathSymbol{\boxtimes}{\mathbin}{AMSa}{"02}

%
% Macros for moduli stacks/spaces
%
\def\QCohstack{\mathcal{QC}\!\mathit{oh}}
\def\Cohstack{\mathcal{C}\!\mathit{oh}}
\def\Spacesstack{\mathcal{S}\!\mathit{paces}}
\def\Quotfunctor{\mathrm{Quot}}
\def\Hilbfunctor{\mathrm{Hilb}}
\def\Curvesstack{\mathcal{C}\!\mathit{urves}}
\def\Polarizedstack{\mathcal{P}\!\mathit{olarized}}
\def\Complexesstack{\mathcal{C}\!\mathit{omplexes}}
% \Pic is the operator that assigns to X its picard group, usage \Pic(X)
% \Picardstack_{X/B} denotes the Picard stack of X over B
% \Picardfunctor_{X/B} denotes the Picard functor of X over B
\def\Pic{\mathop{\mathrm{Pic}}\nolimits}
\def\Picardstack{\mathcal{P}\!\mathit{ic}}
\def\Picardfunctor{\mathrm{Pic}}
\def\Deformationcategory{\mathcal{D}\!\mathit{ef}}
\usepackage{cleveref}
\crefname{originallemma}{original lemma}{original lemmas}
\Crefname{originallemma}{Original lemma}{Original lemmas}
\makeatletter
\newlabel{lemma-coherent-Noetherian}{{07UB}{source}{}{sourcepointer-1}{}}
\newlabel{lemma-coherent-Noetherian@cref}{{[originallemma][0][]07UB}{[1][0][]source}{}{}{}}
\makeatother
\crefname{originaldefinition}{original definition}{original definitions}
\Crefname{originaldefinition}{Original definition}{Original definitions}
\makeatletter
\newlabel{spaces-morphisms-definition-scheme-theoretic-support}{{07U1}{source}{}{sourcepointer-2}{}}
\newlabel{spaces-morphisms-definition-scheme-theoretic-support@cref}{{[originaldefinition][0][]07U1}{[1][0][]source}{}{}{}}
\makeatother
\crefname{originallemma}{original lemma}{original lemmas}
\Crefname{originallemma}{Original lemma}{Original lemmas}
\makeatletter
\newlabel{spaces-morphisms-lemma-immersion-locally-finite-type}{{06ED}{source}{}{sourcepointer-3}{}}
\newlabel{spaces-morphisms-lemma-immersion-locally-finite-type@cref}{{[originallemma][0][]06ED}{[1][0][]source}{}{}{}}
\makeatother
\crefname{originallemma}{original lemma}{original lemmas}
\Crefname{originallemma}{Original lemma}{Original lemmas}
\makeatletter
\newlabel{spaces-morphisms-lemma-locally-finite-type-locally-noetherian}{{04ZK}{source}{}{sourcepointer-4}{}}
\newlabel{spaces-morphisms-lemma-locally-finite-type-locally-noetherian@cref}{{[originallemma][0][]04ZK}{[1][0][]source}{}{}{}}
\makeatother
\title{The module on the scheme theoretic support of a coherent sheaf on a locally Noetherian algebraic space is coherent}
\author{The Stacks Project Authors}
\date{stacks-project:spaces-cohomology.tex:local-snapshot-20261012; GNU FDL 1.2 or later}
\begin{document}
\maketitle
\noindent\textit{Editorial note.} All delivered mathematical text is the author's own, copied from the pinned Stacks Project chapter file spaces-cohomology.tex; no statement or proof step was rewritten, reordered or removed.\par\medskip
\noindent\textit{Editorial note.} Preamble provenance: the pinned chapter file opens with \textbackslash{}input\{preamble\} and carries no preamble text of its own, so the project's single Stacks-wide preamble staged at sources/stacks-spaces-cohomology/preamble.tex is excerpted verbatim as three spans, lines 8-12 (verbatim package and the project's own reference/slogan/history display environments), lines 14-18 (Xy-pic setup) and lines 154-217 (theorem declarations and mathematical macros). The wrapper class and the amsmath/amssymb/amsthm/graphicx base packages are supplied by the assembler. No macro was invented or redefined.\par\medskip
\noindent\textit{Editorial note.} Same-chapter pointers are located at their \textbackslash{}label lines inside the pinned chapter file; cross-chapter pointers are located in the project's own tag list staged in the snapshot (sources/stacks-chow/tags.txt), at the verbatim lines carrying their canonical tags.\par\medskip
\noindent\textit{Allowed standard prerequisites (editorial).}\begin{itemize}
\item The scheme theoretic support of a finite type module, and the module it carries. Morphisms of Spaces, definition spaces-morphisms-definition-scheme-theoretic-support (cited at statement line 2241). It provides the closed immersion \$i : Z \textbackslash{}to X\$ and the quasi-coherent module \$\textbackslash{}mathcal\{G\}\$ with \$i\_*\textbackslash{}mathcal\{G\} = \textbackslash{}mathcal\{F\}\$, and makes the statement meaningful because a coherent module is of finite type.
\item A closed immersion is locally of finite type, and a space locally of finite type over a locally Noetherian space is locally Noetherian. Morphisms of Spaces, labels spaces-morphisms-lemma-immersion-locally-finite-type and spaces-morphisms-lemma-locally-finite-type-locally-noetherian (invoked at proof lines 2251 and 2252). They show that the scheme theoretic support \$Z\$ is itself locally Noetherian, which is the hypothesis under which coherence of \$\textbackslash{}mathcal\{G\}\$ can be tested.
\item Coherence on a locally Noetherian space equals finite type. Same chapter spaces-cohomology.tex, label lemma-coherent-Noetherian (invoked at proof line 2255). Since \$\textbackslash{}mathcal\{G\}\$ is of finite type, this gives that it is coherent.
\end{itemize}\medskip
\ifdefined\hypertarget\hypertarget{sourcepointer-1}{}\fi
\noindent\textit{Source reference.} Original lemma 07UB: sources/stacks-spaces-cohomology/pinned-original-source.tex, source lines 2094--2094. Navigation pointer for \textbackslash{}ref\{lemma-coherent-Noetherian\} used inside the counted span (line 2255); the argument that consumes it is retained in the span.\par\medskip
\ifdefined\hypertarget\hypertarget{sourcepointer-2}{}\fi
\noindent\textit{Source reference.} Original definition 07U1: sources/stacks-chow/tags.txt, source lines 9956--9956. Navigation pointer for \textbackslash{}ref\{spaces-morphisms-definition-scheme-theoretic-support\} used inside the counted span (line 2242); the argument that consumes it is retained in the span.\par\medskip
\ifdefined\hypertarget\hypertarget{sourcepointer-3}{}\fi
\noindent\textit{Source reference.} Original lemma 06ED: sources/stacks-chow/tags.txt, source lines 8212--8212. Navigation pointer for \textbackslash{}ref\{spaces-morphisms-lemma-immersion-locally-finite-type\} used inside the counted span (line 2251); the argument that consumes it is retained in the span.\par\medskip
\ifdefined\hypertarget\hypertarget{sourcepointer-4}{}\fi
\noindent\textit{Source reference.} Original lemma 04ZK: sources/stacks-chow/tags.txt, source lines 6468--6468. Navigation pointer for \textbackslash{}ref\{spaces-morphisms-lemma-locally-finite-type-locally-noetherian\} used inside the counted span (line 2252); the argument that consumes it is retained in the span.\par\medskip
\begin{lemma}
\label{lemma-coherent-support-closed}
Let $S$ be a scheme. Let $X$ be a locally Noetherian algebraic space over $S$.
Let $\mathcal{F}$ be a coherent $\mathcal{O}_X$-module. Let $i : Z \to X$
be the scheme theoretic support of $\mathcal{F}$ and $\mathcal{G}$
the quasi-coherent $\mathcal{O}_Z$-module such that
$i_*\mathcal{G} = \mathcal{F}$, see
Morphisms of Spaces, Definition
\ref{spaces-morphisms-definition-scheme-theoretic-support}.
Then $\mathcal{G}$ is a coherent $\mathcal{O}_Z$-module.
\end{lemma}

\begin{proof}
The statement of the lemma makes sense as a coherent module is in
particular of finite type. Moreover, as $Z \to X$ is a closed immersion
it is locally of finite type and hence $Z$ is locally Noetherian, see
Morphisms of Spaces, Lemmas
\ref{spaces-morphisms-lemma-immersion-locally-finite-type} and
\ref{spaces-morphisms-lemma-locally-finite-type-locally-noetherian}.
Finally, as $\mathcal{G}$ is of finite type it is a coherent
$\mathcal{O}_Z$-module by
Lemma \ref{lemma-coherent-Noetherian}
\end{proof}

\end{document}
